\chapter{METODOLOGIA}
\label{chap:metodologia}
O presente trabalho se iniciou com uma análise na literatura para verificar as propostas atualmente existentes para geração de consultas GraphQL de forma aleatória. A pesquisa serviu como base para analisar as atuais proposições para tal tema e explorar ideias para o desenvolvimento do projeto. As etapas para a realização da pesquisa encontram-se descritas a seguir:

\begin{itemize}
\item Fase 1: Revisão bibliográfica sobre o tema abordado. Nessa etapa foram analisados diversos trabalhos propostos na área de geração de consultas GraphQL de forma aleatória, bem como na área de testes de APIs WEB.
\item Fase 2: Analisar o formato e especificação da linguagem. Etapa onde se faz necessário analisar a especificação da linguagem para verificar as regras, para criação de consultas válidas que obedecem a sintaxe da ferramenta.
\item Fase 3: Criar ferramenta para geração das consultas. Fase onde, com a especificação da linguagem GraphQL, faz-se a criação de ferramenta para a geração das consultas utilizando a linguagem de programação funcional Haskell.
\item Fase 4: Aplicação dos testes baseados em propriedade. Com as consultas geradas deve-se aplicar os testes baseados em propriedades para verificar se as consultas geradas são válidas.
\item Fase 5: Testar as consultas geradas em uma API. Nesta fase, deve-se realizar testes em uma API GraphQL para verificar se as consultas são interpretadas corretamente quando utilizado em um cenário de teste.
\end{itemize}
